\documentclass[12pt]{article}
\usepackage[margin=1in]{geometry}
\usepackage{amsmath}
\usepackage{amssymb}
\usepackage{booktabs}
\usepackage{hyperref}
\usepackage{listings}
\usepackage{xcolor}
\usepackage{titlesec}
\usepackage{parskip}
\usepackage{graphicx}
\usepackage{float}

\hypersetup{
    colorlinks=true,
    linkcolor=blue,
    urlcolor=blue
}

\lstset{
    basicstyle=\ttfamily\small,
    backgroundcolor=\color{gray!10},
    frame=single,
    breaklines=true,
    keywordstyle=\color{blue},
    commentstyle=\color{green!60!black},
    stringstyle=\color{orange}
}

\title{\textbf{Phase 1 Simulation Report} \\
\large Surface-Proximate Drone Docking Project}
\author{Sean Dai}
\date{January 20th, 2026}

\begin{document}

\maketitle
\tableofcontents
\newpage

% -------------------------------------------------------
\section{Overview}
% -------------------------------------------------------

Phase 1 of the surface-proximate drone docking project established a
complete Python simulation environment for a 2D planar quadrotor operating
near a vertical wall. No hardware was used. All work ran on a laptop using
\texttt{numpy}, \texttt{scipy}, and \texttt{matplotlib}.

The three core goals of Phase 1 were:

\begin{enumerate}
    \item Verify stable free-air hover with position error below 5\,cm.
    \item Confirm that wall disturbance terms produce a visible perturbation
          on the drone trajectory.
    \item Demonstrate that the Dockability Margin $M$ decreases as the drone
          approaches the wall and responds correctly to disturbance strength.
\end{enumerate}

All three goals were achieved by the end of Phase 1.

% -------------------------------------------------------
\section{File Structure and Responsibilities}
% -------------------------------------------------------

The simulation was organized into seven files, each with a single
well-defined responsibility.

\begin{center}
\begin{tabular}{ll}
\toprule
\textbf{File} & \textbf{Responsibility} \\
\midrule
\texttt{params.py}     & All numerical constants in one dictionary \\
\texttt{dynamics.py}   & Equations of motion and wall disturbance model \\
\texttt{controller.py} & Cascaded PD nominal controller \\
\texttt{margin.py}     & Dockability Margin and six sub-margin computations \\
\texttt{logger.py}     & Per-timestep data logging and \texttt{.npz} saving \\
\texttt{run\_sim.py}   & Simulation entry point and experiment configuration \\
\texttt{visualize.py}  & All plotting functions \\
\bottomrule
\end{tabular}
\end{center}

% -------------------------------------------------------
\section{System Model}
% -------------------------------------------------------

\subsection{State Representation}

The drone is modeled as a rigid body in the vertical plane. The six-element
state vector is:

\begin{equation}
    \mathbf{x} = \begin{bmatrix} x & z & v_x & v_z & \theta & \dot{\theta} \end{bmatrix}^\top
\end{equation}

where $x$ is horizontal position, $z$ is altitude, $v_x$ and $v_z$ are
translational velocities, $\theta$ is pitch angle, and $\dot{\theta}$ is
pitch rate.

\subsection{Equations of Motion}

The translational and rotational dynamics are:

\begin{align}
    m\dot{v}_x &= -T_\mathrm{eff}\sin\theta + f_\mathrm{wall}(d, v_n) \\
    m\dot{v}_z &= T_\mathrm{eff}\cos\theta - mg \\
    I\ddot{\theta} &= \tau
\end{align}

where the wall-modified effective thrust is:

\begin{equation}
    T_\mathrm{eff}(d) = T \cdot \Gamma(d), \qquad
    \Gamma(d) = 1 + \alpha e^{-\beta d}
\end{equation}

and the wall disturbance force is:

\begin{equation}
    f_\mathrm{wall}(d, v_n) = -\left(c_0 e^{-\lambda d}
    + c_1 e^{-\lambda d} v_n
    + c_2 e^{-\lambda d} v_n |v_n|\right)
\end{equation}

Here $d = x_w - x$ is the signed distance to the wall and
$v_n = -v_x$ is the wall-normal velocity (positive = receding).

% -------------------------------------------------------
\section{Controller Design}
% -------------------------------------------------------

\subsection{Cascaded PD Architecture}

A two-level cascaded PD controller was implemented. The outer loop
converts position errors into desired pitch angle and thrust; the inner
loop converts attitude error into torque.

\textbf{Outer loop --- vertical channel:}
\begin{align}
    a_{z,\mathrm{des}} &= k_{p,z}(z_\mathrm{des} - z) + k_{d,z}(0 - v_z) \\
    T &= m(g + a_{z,\mathrm{des}}), \quad T \in [0, T_\mathrm{max}]
\end{align}

\textbf{Outer loop --- horizontal channel:}
\begin{align}
    a_{x,\mathrm{des}} &= k_{p,x}(x_\mathrm{des} - x) + k_{d,x}(0 - v_x) \\
    \theta_\mathrm{des} &= -\frac{m \cdot a_{x,\mathrm{des}}}{T},
    \quad \theta_\mathrm{des} \in [-\theta_\mathrm{max}, +\theta_\mathrm{max}]
\end{align}

\textbf{Inner loop --- attitude channel:}
\begin{equation}
    \tau = I\left(k_{p,\theta}(\theta_\mathrm{des} - \theta)
    + k_{d,\theta}(0 - \dot{\theta})\right)
\end{equation}

\subsection{Final Tuned Gains}

Significant gain tuning was required to achieve the 5\,cm hover criterion.
The final values used are shown below.

\begin{center}
\begin{tabular}{lll}
\toprule
\textbf{Parameter} & \textbf{Final Value} & \textbf{Description} \\
\midrule
$k_{p,z}$ & 8.0  & Vertical position P gain \\
$k_{d,z}$ & 4.0  & Vertical velocity D gain \\
$k_{p,x}$ & 2.0  & Horizontal position P gain \\
$k_{d,x}$ & 10.0 & Horizontal velocity D gain \\
$k_{p,\theta}$ & 15.0 & Pitch angle P gain \\
$k_{d,\theta}$ & 15.0 & Pitch rate D gain \\
\bottomrule
\end{tabular}
\end{center}

% -------------------------------------------------------
\section{Dockability Margin}
% -------------------------------------------------------

\subsection{Sub-Margin Definitions}

The Dockability Margin $M \in [0,1]$ is the minimum of six normalized
sub-margins, each computed via the saturation function
$\mathrm{sat}(x) = \min(1, \max(0, x))$.

\begin{center}
\begin{tabular}{lll}
\toprule
\textbf{Sub-margin} & \textbf{Formula} & \textbf{Measures} \\
\midrule
$M_d$ & $\mathrm{sat}\!\left(\dfrac{d - d_{\min}}{d_\mathrm{safe} - d_{\min}}\right)$ & Wall clearance \\[8pt]
$M_v$ & $\mathrm{sat}\!\left(1 - \dfrac{\max(0,-v_n)}{v_{n,\max}}\right)$ & Approach speed \\[8pt]
$M_\delta$ & $\mathrm{sat}\!\left(1 - \dfrac{|\theta|}{\delta_{\max}}\right)$ & Tilt alignment \\[8pt]
$M_\omega$ & $\mathrm{sat}\!\left(1 - \dfrac{|\dot{\theta}|}{\omega_{\max}}\right)$ & Angular rate \\[8pt]
$M_u$ & $\mathrm{sat}\!\left(1 - \dfrac{T_\mathrm{cmd}}{T_{\max}}\right)$ & Control authority \\[8pt]
$M_r$ & $\mathrm{sat}\!\left(1 - \dfrac{|T_\mathrm{eff} - T_\mathrm{cmd}|/m}{r_{\max}}\right)$ & Disturbance residual \\
\bottomrule
\end{tabular}
\end{center}

\begin{equation}
    M = \min\{M_d,\, M_v,\, M_\delta,\, M_\omega,\, M_u,\, M_r\}
\end{equation}

\subsection{Calibrated Thresholds}

Based on Experiment 3 results, the following thresholds were calibrated:

\begin{center}
\begin{tabular}{lll}
\toprule
\textbf{Threshold} & \textbf{Value} & \textbf{Meaning} \\
\midrule
$M_\mathrm{safe}$  & 0.6  & CBF layer activates below this value \\
$M_\mathrm{dock}$  & 0.4  & Docking mode entry \\
$M_\mathrm{abort}$ & 0.15 & Abort and retreat triggered \\
\bottomrule
\end{tabular}
\end{center}

% -------------------------------------------------------
\section{Experiments and Results}
% -------------------------------------------------------

\subsection{Experiment 1 --- Free-Air Hover}

\textbf{Purpose:} Verify the controller and dynamics work correctly before
adding wall effects.

\textbf{Setup:} Target $(x, z) = (1.0, 1.5)$\,m, initial state
$(1.2, 1.5)$\,m, \texttt{wall\_on=False}, duration 15\,s.

\textbf{Result: PASS.}
\begin{itemize}
    \item $x$ error entered the 5\,cm threshold at approximately $t = 4$\,s.
    \item $z$ error remained below 5\,cm for the entire run
          (sub-millimeter steady state).
    \item Thrust at steady-state hover $\approx 9.81$\,N $= mg$.
    \item No sustained oscillation observed after $t = 4$\,s.
\end{itemize}

\subsection{Experiment 2 --- Near-Wall Hover}

\textbf{Purpose:} Verify that wall disturbance terms produce a visible
perturbation and that the controller responds correctly.

\textbf{Setup:} Target $(x, z) = (1.8, 1.5)$\,m, \texttt{wall\_on=True},
duration 20\,s. Wall at $x_w = 2.0$\,m, giving $d \approx 0.2$\,m.

\textbf{Result: PASS.}
\begin{itemize}
    \item $f_\mathrm{wall}$ was negative throughout, confirming wall
          attraction force.
    \item $f_\mathrm{wall}$ magnitude increased as $d$ decreased, ranging
          from $-0.30$\,N at $d = 0.45$\,m to $-0.50$\,N at $d = 0.28$\,m.
    \item $\Gamma > 1.0$ everywhere in the range, reaching $\approx 1.10$
          at closest approach.
    \item Drone trajectory showed visible oscillation caused by the
          disturbance force, confirming the wall effect is active.
\end{itemize}

\subsection{Experiment 3 --- Constant-Speed Wall Approach}

\textbf{Purpose:} Verify that $M$ responds correctly during a controlled
wall approach.

\textbf{Setup:} Moving target $x_\mathrm{des}(t) = \min(x_\mathrm{stop},\,
1.5 + 0.05t)$, $z_\mathrm{des} = 1.5$\,m, \texttt{wall\_on=True},
duration 25\,s.

\textbf{Result: PASS.}
\begin{itemize}
    \item $M$ showed a clear general downward trend from $\approx 0.7$
          at $t = 0$ to $\approx 0.15$ near the wall.
    \item $M_d$ (clearance) was the dominant sub-margin driving the decrease.
    \item $M$ crossed all three threshold lines ($M_\mathrm{safe}$,
          $M_\mathrm{dock}$, $M_\mathrm{abort}$) during the approach.
\end{itemize}

\subsection{Experiment 4 --- Sensitivity to Disturbance Strength}

\textbf{Purpose:} Verify that $M$ responds correctly to different
disturbance levels and that $M_r$ is most sensitive to disturbance strength.

\textbf{Setup:} Three runs with the same approach profile but different
disturbance coefficients:

\begin{center}
\begin{tabular}{llll}
\toprule
\textbf{Run} & $c_0$ & $c_1$ & $c_2$ \\
\midrule
No disturbance    & 0.0 & 0.0 & 0.0 \\
Default           & 0.8 & 0.3 & 0.1 \\
Double strength   & 1.6 & 0.6 & 0.2 \\
\bottomrule
\end{tabular}
\end{center}

\textbf{Result: PASS.}
\begin{itemize}
    \item Default disturbance produced the smoothest and most monotonic
          decrease in $M$, reaching $\approx 0.05$ at end of run.
    \item Double disturbance caused large oscillations in $M$ between
          $0.3$ and $0.75$, demonstrating that stronger disturbances
          destabilize the margin.
    \item No disturbance (green curve) overlapped with default at the
          start, confirming the margin system correctly distinguishes
          disturbance levels.
\end{itemize}

% -------------------------------------------------------
\section{Challenges Encountered}
% -------------------------------------------------------

\subsection{Sign Error in Horizontal Controller}

The most impactful bug discovered during Phase 1 was a sign error in the
horizontal channel of the controller. The line:

\begin{lstlisting}[language=Python]
theta_des = (m * ax_des) / T
\end{lstlisting}

should have been:

\begin{lstlisting}[language=Python]
theta_des = -(m * ax_des) / T
\end{lstlisting}

The dynamics equation $m\dot{v}_x = -T\sin\theta$ requires a negative sign
when converting desired acceleration to desired pitch angle. Without this
fix, the drone drifted to over 450\,m horizontally instead of converging
to the target.

\subsection{Controller Gain Tuning}

Achieving the 5\,cm hover criterion within 8\,s required extensive gain
tuning across more than 15 simulation runs. The key challenge was the
competing requirements of the cascaded loops:

\begin{itemize}
    \item Higher $k_{p,x}$ caused oscillations because the outer loop
          competed with the inner loop at similar timescales.
    \item Higher $k_{d,x}$ alone was insufficient to suppress oscillations
          without also increasing $k_{d,\theta}$.
    \item The final solution required making $k_{d,\theta} = 15.0$
          significantly larger than its original value of 6.0, which
          effectively decoupled the two loops by making the inner loop
          much faster.
\end{itemize}

\subsection{Experiment 3 Approach Direction Bug}

The initial Experiment 3 implementation used:

\begin{lstlisting}[language=Python]
x_des = lambda t: max(d_min + 0.02, x_start - v_approach * t)
\end{lstlisting}

which moved the target \emph{away} from the wall (decreasing $x$) instead
of toward it. The fix was to change the sign and swap \texttt{max} for
\texttt{min}:

\begin{lstlisting}[language=Python]
x_des = lambda t: min(x_stop, x_start + v_approach * t)
\end{lstlisting}

\subsection{Python Module Caching}

Several debugging sessions were extended because changes to source files
were not reflected in the running Python session due to module caching.
The fix was always to either restart the Python interpreter or use
\texttt{importlib.reload()} and to delete \texttt{\_\_pycache\_\_} when
needed.

\subsection{Dockability Margin Starting Below 1.0}

The overall margin $M$ never started near 1.0 as expected because
$M_u$ (control authority) was capped at $\approx 0.6$ even at rest.
This is because the drone uses approximately half its available thrust
just to hover against gravity. The fix was to double $T_\mathrm{max}$
from 19.62\,N to 39.24\,N, giving the controller more headroom and
raising $M_u$ to $\approx 0.75$ at hover.

% -------------------------------------------------------
\section{Final Parameter Values}
% -------------------------------------------------------

\begin{center}
\begin{tabular}{lll}
\toprule
\textbf{Parameter} & \textbf{Final Value} & \textbf{Description} \\
\midrule
$m$               & 1.0   & Vehicle mass (kg) \\
$I$               & 0.01  & Pitch inertia (kg\,m$^2$) \\
$g$               & 9.81  & Gravity (m/s$^2$) \\
$T_\mathrm{max}$  & 39.24 & Max thrust (N) \\
$x_w$             & 2.0   & Wall position (m) \\
$\alpha$          & 0.15  & Thrust modification strength \\
$\beta$           & 1.5   & Thrust modification decay rate (m$^{-1}$) \\
$\lambda$         & 2.0   & Disturbance force decay rate (m$^{-1}$) \\
$c_0$             & 0.8   & Pressure attraction coefficient (N) \\
$c_1$             & 0.3   & Viscous damping coefficient (N\,s/m) \\
$c_2$             & 0.1   & Nonlinear drag coefficient (N\,s$^2$/m$^2$) \\
$d_{\min}$          & 0.05  & Minimum wall clearance (m) \\
$d_\mathrm{safe}$ & 0.80  & Nominal safe distance (m) \\
$v_{n,\max}$      & 1.5   & Max allowable approach speed (m/s) \\
$\delta_{\max}$     & 0.349 & Max alignment error (rad) \\
$\omega_{\max}$     & 2.0   & Max angular rate (rad/s) \\
$r_{\max}$          & 2.0   & Disturbance residual threshold (m/s$^2$) \\
\bottomrule
\end{tabular}
\end{center}

% -------------------------------------------------------
\section{Notes for Phase 2}
% -------------------------------------------------------

\subsection{CBF Implementation}

Phase 2 requires installing \texttt{cvxpy} for the quadratic program solver:

\begin{lstlisting}[language=bash]
pip install cvxpy
\end{lstlisting}

The CBF safety filter will receive the nominal control input
$(T_\mathrm{nom}, \tau_\mathrm{nom})$ from \texttt{controller.py} and solve:

\begin{equation}
    \min_u \|u - u_\mathrm{nom}\|^2 \quad \text{subject to} \quad
    \dot{h}_i(x, u) + \alpha_i(h_i(x)) \geq 0 \quad \forall i
\end{equation}

Barrier functions are needed for minimum wall clearance, maximum approach
velocity, attitude limits, and actuator saturation.

\subsection{State Machine Thresholds}

The M thresholds calibrated in Phase 1 that will govern state machine
transitions in Phase 2 are:

\begin{center}
\begin{tabular}{lll}
\toprule
\textbf{Threshold} & \textbf{Value} & \textbf{Trigger} \\
\midrule
$M_\mathrm{safe}$  & 0.6  & Enter surface-proximate mode, CBF activates \\
$M_\mathrm{dock}$  & 0.4  & Enter docking mode \\
$M_\mathrm{abort}$ & 0.15 & Abort and retreat from wall \\
\bottomrule
\end{tabular}
\end{center}

\subsection{Known Limitations to Address}

\begin{itemize}
    \item The controller oscillates significantly in the presence of strong
          wall disturbances. The CBF layer should suppress unsafe approach
          velocities that cause these oscillations.
    \item $M$ is not perfectly monotonic during approach due to controller
          oscillations. Phase 2 should produce smoother M profiles because
          the CBF limits approach velocity.
    \item The current controller does not reach $d_{\min}$ before the
          simulation ends. The docking mode in Phase 2 should handle the
          final close-range approach more carefully.
    \item The 2D planar model will need to be extended toward a 3D model
          in later phases for hardware validation.
\end{itemize}

\subsection{Files to Create in Phase 2}

\begin{center}
\begin{tabular}{ll}
\toprule
\textbf{File} & \textbf{Responsibility} \\
\midrule
\texttt{cbf.py}           & CBF safety filter and QP solver \\
\texttt{state\_machine.py} & Hybrid docking state machine \\
\bottomrule
\end{tabular}
\end{center}

Existing files \texttt{run\_sim.py}, \texttt{logger.py}, and
\texttt{visualize.py} will also need updates to incorporate the new
CBF and state machine outputs.

% -------------------------------------------------------
\section{Phase 1 Success Checklist}
% -------------------------------------------------------

\begin{center}
\begin{tabular}{cl}
\toprule
\textbf{Status} & \textbf{Item} \\
\midrule
$\checkmark$ & Experiment 1 runs to completion with no errors \\
$\checkmark$ & Hover position error $<$ 5\,cm in both $x$ and $z$ \\
$\checkmark$ & Thrust at steady-state hover $\approx 9.81$\,N $= mg$ \\
$\checkmark$ & $f_\mathrm{wall}$ is non-zero and negative in Experiment 2 \\
$\checkmark$ & $\Gamma > 1.0$ at $d < 0.5$\,m in Experiment 2 \\
$\checkmark$ & $M$ decreases during Experiment 3 \\
$\checkmark$ & $M_d$ is the dominant sub-margin during approach \\
$\checkmark$ & Experiment 4 shows stronger disturbance $\Rightarrow$ lower $M$ \\
$\checkmark$ & All four experiment logs saved as \texttt{.npz} files \\
$\checkmark$ & Plots saved to \texttt{plots/} directory \\
\bottomrule
\end{tabular}
\end{center}

\end{document}
